\section{Why standard balancing and capacity do not force flatness}
\label{sec:capacity-obstruction}

Operator scaling connects invariant theory, noncommutative rank and
capacity; primary sources include Garg--Gurvits--Oliveira--Wigderson,
\emph{Operator scaling: theory and applications},
\url{https://arxiv.org/abs/1511.03730}, and Franks,
\emph{Operator scaling with specified marginals},
\url{https://arxiv.org/abs/1801.01412}. These theorems suggest a natural
test: can balancing the entry magnitudes of an orthogonal matrix force
a flat representative? The following elementary obstructions identify
precisely what that test loses. The arguments below are independently
derived; no external implementation is used.

For $U\in O(n)$ write $B_{ij}=u_{ij}^2$. Then $B$ is doubly stochastic:
it is nonnegative and every row and column sums to one. A matrix of
this form is called orthostochastic. Flatness asks for the very special
orthostochastic matrix $J_n/n$, not merely for prescribed marginals.
The distance-minimization formulation itself is established literature;
see Arasu--Mohan, \emph{Entropy of orthogonal matrices and minimum
distance orthostochastic matrices from the uniform van der Waerden
matrices}, Discrete Optimization 31 (2019), 115--144,
\url{https://doi.org/10.1016/j.disopt.2018.10.001}.

\begin{proposition}[Rigidity of positive diagonal balancing]
Let $U\in O(n)$ and let $D,E$ be positive diagonal matrices.
If $DUE\in O(n)$, then $DUE=U$ entry by entry.
If the bipartite support graph of $U$ is connected, necessarily
$D=cI_n$ and $E=c^{-1}I_n$ for some $c>0$.
\end{proposition}
\begin{proof}
Put $a_i=D_{ii}^2$, $b_j=E_{jj}^2$, and
$B'_{ij}=a_iB_{ij}b_j$. Both $B,B'$ are doubly stochastic and have
the same support. For equally normalized nonnegative arrays, relative
entropy is nonnegative, with equality exactly when the arrays agree.
Here direct summation using the marginals gives
\[
 \sum_{ij:B_{ij}>0}B'_{ij}\log\frac{B'_{ij}}{B_{ij}}
 =\sum_i\log a_i+\sum_j\log b_j,
\]
whereas the reverse relative entropy is the negative of that expression.
Both are nonnegative, so both vanish and $B'=B$.
Positive diagonal multipliers preserve the signs of every nonzero
entry, proving $DUE=U$. On each support edge $a_ib_j=1$;
connectedness propagates a common value of $a_i$ and its reciprocal
among the $b_j$, giving the final statement.
\end{proof}

This does not claim that general left-right matrix scaling is powerless.
It says that positive \emph{diagonal} scaling constrained to preserve
orthogonality cannot change the matrix at all.

\begin{proposition}[The natural classical and operator capacities are blind]
Let $B$ be any doubly stochastic matrix. Define
\[
 p_B(x)=\prod_{i=1}^n\left(\sum_j B_{ij}x_j\right),\qquad
 \operatorname{cap}(p_B)=\inf_{x_j>0}\frac{p_B(x)}{\prod_jx_j}.
\]
Also define the completely positive map
\[
 T_B(X)=\operatorname{diag}\bigl(B\operatorname{diag}(X)\bigr),
 \qquad
 \operatorname{cap}(T_B)=\inf_{X\succ0}
       \frac{\det T_B(X)}{\det X}.
\]
Then both capacities equal $1$. Moreover
$T_B(I)=T_B^*(I)=I$. Thus every orthogonal $U$, flat or not,
already passes these balancing and capacity tests for $B=U\circ U$.
\end{proposition}
\begin{proof}
Weighted arithmetic--geometric mean, followed by the column sums, gives
\[
 \prod_i\sum_jB_{ij}x_j
 \ge\prod_{ij}x_j^{B_{ij}}=\prod_jx_j.
\]
Equality is achieved at $x=(1,\ldots,1)$, proving the polynomial claim.
The Kraus matrices $K_{ij}=\sqrt{B_{ij}}E_{ij}$ represent $T_B$, so it
is completely positive. For $X\succ0$, the preceding inequality and
the determinant bound by the product of diagonal entries yield
\[
 \det T_B(X)\ge\prod_jX_{jj}\ge\det X.
\]
Taking $X=I$ proves the operator-capacity assertion. The row and column
sums give the two marginal identities directly.
\end{proof}

The theorem is a limitation of these particular encodings, not a
limitation of every invariant-theoretic construction. A successful
encoding would have to retain constraints absent from row and column
marginals and these constant capacities, such as simultaneous real
sign cancellation between every pair of rows.

\subsection{A quantitative separation despite perfect capacity}
There is a uniform positive separation from flatness in every forbidden
order $n\equiv2\pmod4$, $n>2$, even though the preceding capacities
remain exactly $1$.

\begin{proposition}[Elementary parity gap]
Let $n\equiv2\pmod4$ and $U\in O(n)$. Then
\[
 F_4(U)-1\ge
 \frac{\left(\sqrt{1+\sqrt{2(n-2)/n}}-1\right)^2}{n}
 \ge \frac{2(n-2)}{9n^2}.
\]
Equivalently this lower bounds
$\|U\circ U-J_n/n\|_F^2$, since that squared distance is $F_4(U)-1$.
No optimality of the constants is asserted.
\end{proposition}
\begin{proof}
For any sign matrix $H$ of order $n$, split its rows into two classes
according to the parity of their number of negative entries.
Two rows in the same class have inner product congruent to $2$
modulo $4$, hence of magnitude at least $2$. If one class has $a$
rows, there are $a(a-1)+(n-a)(n-a-1)\ge n(n-2)/2$ ordered pairs of
distinct rows in the same class. Consequently
\[
 \|HH^{\mathsf T}-nI\|_F^2\ge2n(n-2).
\]
Choose $H_{ij}=\operatorname{sign}(u_{ij})$, with either sign allowed
at zero, set $C=H/\sqrt n$, and let $d=\|C-U\|_F$.
Entrywise,
\[
 \left(|u_{ij}|-n^{-1/2}\right)^2
 =\frac{(u_{ij}^2-1/n)^2}{(|u_{ij}|+n^{-1/2})^2}
 \le n(u_{ij}^2-1/n)^2,
\]
so $d^2\le n(F_4(U)-1)$. Writing $C=U+(C-U)$ and using the
triangle inequality and $\|EE^{\mathsf T}\|_F\le\|E\|_F^2$ gives
\[
 L:=\sqrt{2(n-2)/n}\le\|CC^{\mathsf T}-I\|_F\le2d+d^2.
\]
Hence $d\ge\sqrt{1+L}-1$, proving the first bound.
Finally $0\le L<\sqrt2$ and
$(\sqrt{1+L}-1)^2=L^2/(1+\sqrt{1+L})^2\ge L^2/9$.
\end{proof}

For example, at order $6$ every orthogonal matrix passes both capacity
tests with value exactly $1$, yet $F_4(U)\ge1+2/81$.
At admissible orders the same capacity tests stay constant, so they
cannot by themselves distinguish existence from nonexistence of a flat
point. Proving a stronger real-sign-sensitive feasibility theorem remains
an open research task; replacing it by ordinary bistochastic feasibility
would remove the essential difficulty.
